\documentclass{leveldoc}
\DocUses{foundations,slingshot,first_coil,stray_field,mains_hum}
\DocFingerprint{97778a70206f254e}

\begin{document}

\section{What the level asks}\label{sec:asks}

The game's text:
\leveldescription

An electron, charge $q = \val{q}$ and mass $m = \val{m}$, is launched from
$(\val{xZero}, \val{yZero})$ along $+x$ with $T_0 = \val{Tzero}$, in Newton's mechanics
($c = \infty$), and must enter the detector $\val{xDetLo} \le x \le \val{xDetHi}$,
$\val{yDetLo} \le y \le \val{yDetHi}$ under each of the level's three disturbances: no field, and a
uniform stray field $B_z = \pm\val{B}$ perpendicular to the board, ``\val{distNorth}'' and
``\val{distSouth}'' (\physref{2.3}; every shot is flown under every disturbance, and the level is
solved only if every flight arrives). The player places up to four charges from the listed
magnitudes, either sign, at nodes with $\val{xRegLo} \le x \le \val{xRegHi}$,
$\val{yRegLo} \le y \le \val{yRegHi}$. A flight ends when the electron enters the detector, touches
a charge (radius $\val{R}$), or leaves the world (\xsecref{foundations}{sec:board}).

The reference places two negative charges, which repel the electron: $Q = \val{QOne}$ at
$(\val{xOne}, \val{yOne})$, below the beam line near the launch, and $Q = \val{QTwo}$ at
$(\val{xTwo}, \val{yTwo})$, above it further on. The document solves the empty board in closed
form, marches the reference's three flights with the energy as the check, and measures how the
charges make the arrival nearly independent of the stray field: not by cancelling it, but by
folding the flights it bends upward back down.

\section{Setup}\label{sec:setup}

The level file, field by field (the appendix reproduces it):
\setuptable

The electron's speed is $v_0 = \val{vZero}$ and its momentum $p_\mathrm{i} = mv_0$. In the stray
field it circles with the radius $R_B = p_\mathrm{i}/(\abs{q}B) = \val{Rlarmor}$; the magnetic force
on it, $\abs{q}v_0B = \val{forceB}$, is a few hundredths of the reference charges' forces near them
($\abs{qQ}/r^2$ with $qQ = \val{kOne}$ and $\val{kTwo}$). Over the $L = \val{L}$ cells to the
detector a circle of radius $R_B$ departs from its tangent by $R_B - \sqrt{R_B^2 - L^2} = \val{sag}$
($\approx L^2/2R_B = \val{sagApprox}$), against the detector's half-height $\val{halfHeight}$.

\begin{figure}[htbp]
\centering
\begin{tikzpicture}
\begin{axis}[axis equal image, width=0.92\linewidth, xmin=-1.6, xmax=31.6, ymin=-1.6, ymax=21.6,
  xtick={0,5,...,30}, ytick={0,5,...,20}, tick label style={font=\small},
  xlabel={$x$}, ylabel={$y$}, label style={font=\small}]
  \draw[black!8] (0,0) grid (30,20);
  \draw[black!45, dashed] (-1,-1) rectangle (\rawval{xWorld},\rawval{yWorld});
  \draw[black!30, dotted] (\rawval{xRegLo},\rawval{yRegLo}) rectangle (\rawval{xRegHi},\rawval{yRegHi});
  \fill[green!50!black, fill opacity=0.18] (\rawval{xDetLo},\rawval{yDetLo}) rectangle (\rawval{xDetHi},\rawval{yDetHi});
  \draw[green!40!black] (\rawval{xDetLo},\rawval{yDetLo}) rectangle (\rawval{xDetHi},\rawval{yDetHi});
  \addplot[thick, black!70] table {pathoff.dat};
  \addplot[thick, blue!70!black] table {pathnorth.dat};
  \addplot[thick, red!70!black] table {pathsouth.dat};
  \fill[blue!60!black] (\rawval{xOne},\rawval{yOne}) circle[radius=\rawval{R}];
  \fill[blue!60!black] (\rawval{xTwo},\rawval{yTwo}) circle[radius=\rawval{R}];
  \fill (\rawval{xZero},\rawval{yZero}) circle[radius=0.15];
\end{axis}
\end{tikzpicture}
\caption{The reference's negative charges (blue discs), the placement region (dotted) and the
three flights: no field (black), facing north (blue) and facing south (red).}
\label{fig:board}
\end{figure}

\section{Principles and model}\label{sec:model}

\begin{principles}
  \item Newton's equation with the Lorentz force \xeqref{foundations}{eq:eom} ($c = \infty$),
    the charges' Coulomb fields \xeqref{foundations}{eq:coulomb} and the uniform $B_z$
    (\physref{2.3}).
  \item In a static field the energy $p^2/2m + q\phi$ is conserved, the magnetic force doing no
    work (\xsecref{foundations}{sec:energy}); a charge in a uniform $B$ alone moves on a circle,
    \xeqref{first_coil}{eq:larmor-radius}.
  \item The board's rules; the classical Runge--Kutta method, cubic Hermite interpolation and the
    error estimate from the step $2h$; bisection and golden-section search
    (\xsecref{foundations}{sec:board}, \xsecref{foundations}{sec:numerics}), as in
    \xsecref{slingshot}{sec:method}.
\end{principles}

The electron starts at $(\val{xZero}, \val{yZero})$ with $\vect{p} = (p_\mathrm{i}, 0)$ and moves by
\begin{equation}\label{eq:eom-b}
  \dot{\vect{p}} = q\big(\vect{E}_Q(\vect{x}) + \vect{v}\times B_z\uvec{z}\big), \qquad
  \text{in the plane}\quad \dot p_x = q(E_{Q,x} + v_yB_z), \quad \dot p_y = q(E_{Q,y} - v_xB_z),
\end{equation}
until the first of the board's events; the energy
\begin{equation}\label{eq:energy-b}
  W = \frac{p^2}{2m} + \sum_Q \frac{qQ}{\abs{\vect{x} - \vect{r}_Q}}
\end{equation}
is conserved in each of the three fields.

\begin{assumption}{A uniform stray field}\label{as:uniform}
The stray field is uniform and static, an exact solution of Maxwell's equations (\physref{2.3}),
so nothing is neglected within the game's model. The Earth's field is uniform over a lab and is
about $50\,\mu$T; the level's $B$ is chosen to bend the beam by a few cells over the board.
\end{assumption}

\begin{assumption}{Fixed charges, a point electron, no radiation}\label{as:fixed}
The charges stay put and the electron is a point; with $c = \infty$ Liénard's power
\xeqref{foundations}{eq:larmor} vanishes. These are the game's model.
\end{assumption}

\section{Derivation}\label{sec:derivation}

\subsection{The empty board}\label{sec:circle}

Without charges and without field the electron flies straight along $y = \val{yZero}$ into the
detector's left face at $t = \val{freeTOff}$. With $B_z \ne 0$ it moves on the circle of radius
$R_B$ tangent to the launch line; the force $-qv_xB_z$ along $y$ is upward for $-qB_z > 0$, so the
circle curves up facing north ($B_z > 0$, $q < 0$) and down facing south:
\begin{equation}\label{eq:circle}
  y(x) = y_0 \pm \Big(R_B - \sqrt{R_B^2 - (x - x_0)^2}\Big), \qquad
  t(x) = \frac{R_B}{v_0}\arcsin\frac{x - x_0}{R_B}, \qquad
  p_y = \pm p_\mathrm{i}\frac{x - x_0}{R_B}.
\end{equation}
At the detector's face it is $\val{sag}$ off the line, outside the detector; it leaves the world's
right edge at $y = \val{freeYNorth}$ and $\val{freeYSouth}$, at $t = \val{freeTNorth}$. To first
order in $B$ the height at the detector moves by
\begin{equation}\label{eq:sag}
  \deriv{y}{B_z} = \frac{\abs{q}L^2}{2p_\mathrm{i}} = \val{slopeFree}
\end{equation}
per unit of field, and the flights facing north and south end $\val{spreadFree}$ apart.

\subsection{The march}\label{sec:march}

With the charges there is no closed form. \eqref{eq:eom-b} is marched by
\xeqref{foundations}{eq:rk4} at $h = \val{H}$ in each field; the events are located on the cubic
Hermite interpolant by bisection, the error is estimated against the step $2h$, and
\eqref{eq:energy-b} checks every step (\xsecref{stray_field}{sec:march}).

\subsection{How the charges handle the stray field}\label{sec:fold}

Marching flights with the reference's charges in fields between $-1.5B$ and $1.5B$ and recording
where they cross the line $x = \val{xDetLo}$ gives Figure~\ref{fig:curve}. Its slope at
$B_z = 0$, by central differences, is $\val{slopeRef}$, a factor $\val{slopeRatio}$ smaller than the
empty board's \eqref{eq:sag}; but the curve is far from straight. Without field the electron
crosses the line at $\val{yLineZero}$, above the detector: the lower charge, which it passes first,
has lifted it, and it comes down into the detector's top face further on. Facing north the field
bends it up, towards the upper charge (it passes $\val{nearNorth}$ from it, against $\val{nearOff}$
without field), whose repulsion grows faster than the lift: the crossing height is highest,
$\val{foldY}$, at $B_z \approx \val{foldB}$ and falls beyond, to $\val{yLineNorth}$ at $B_z = \val{B}$.
Facing south the field bends the electron down and away from that charge (to $\val{nearSouth}$),
and it crosses at $\val{yLineSouth}$. The three flights cross within $\val{spreadRef}$ of each other
in height instead of $\val{spreadFree}$: the upper charge folds the upward bend back, as the lower
charge of \levelref{mains_hum} folds its downward tilts.

\begin{figure}[htbp]
\centering
\begin{tikzpicture}
\begin{axis}[width=0.72\linewidth, height=5.4cm, xmin=-11600, xmax=11600, ymin=4, ymax=16,
  xlabel={stray field $B_z$}, ylabel={height at $x = \val{xDetLo}$},
  tick label style={font=\small}, label style={font=\small},
  xtick={-10000,-5000,0,5000,10000}, scaled x ticks=false,
  xticklabel style={/pgf/number format/fixed}]
  \fill[green!50!black, fill opacity=0.15] (-11600,\rawval{yDetLo}) rectangle (11600,\rawval{yDetHi});
  \addplot[thick, blue!70!black] table {curve.dat};
  \addplot[thin, dashed, black!60] table {curvefree.dat};
\end{axis}
\end{tikzpicture}
\caption{Where a flight crosses the line $x = \val{xDetLo}$, against the stray field: with the
reference's charges (solid) and on the empty board (dashed, \eqref{eq:circle}). Shaded: the
detector's span; the level's fields are $0$ and $\pm\val{B}$.}
\label{fig:curve}
\end{figure}

\section{The reference solution by hand}\label{sec:hand}

\begin{procedure}
  \item\label{st:off} \emph{No field.} The march of \S\ref{sec:march} with $B_z = 0$: the electron
    enters the detector through its \val{faceOff} face at $(\val{xOff}, \val{yOff})$ at
    $t = \val{tOff}$, with $\vect{p} = (\val{pxOff}, \val{pyOff})$ (step estimate $\val{errOff}$; the
    energy drifts by $\val{wDevOff}$).
  \item\label{st:north} \emph{Facing north}, $B_z = \val{B}$: through its \val{faceNorth} face at
    $(\val{xNorth}, \val{yNorth})$ at $t = \val{tNorth}$, $\vect{p} = (\val{pxNorth}, \val{pyNorth})$
    (step estimate $\val{errNorth}$; energy drift $\val{wDevNorth}$).
  \item\label{st:south} \emph{Facing south}, $B_z = -\val{B}$: through its \val{faceSouth} face at
    $(\val{xSouth}, \val{ySouth})$ at $t = \val{tSouth}$, $\vect{p} = (\val{pxSouth}, \val{pySouth})$
    (step estimate $\val{errSouth}$; energy drift $\val{wDevSouth}$).
\end{procedure}

\begin{result}[The reference flights]
In all three fields the electron enters the detector. The level is solved.
\end{result}

\section{Comparison with the game}\label{sec:comparison}

\comparisontable

All differences are at most $\val{largest}$, the game's verified accuracy (Foundations §6). The
tolerances are ten times the larger of the two computations' estimates, or the game's own
preview--verification gap where that is larger (the field-free flight, which enters through the
top face, where its point of entry moves along the face by the height's error divided by its
shallow slope); the empty board's values come from the closed form \eqref{eq:circle}.

\section{Other solutions}\label{sec:others}

The empty board fails facing north and south (\S\ref{sec:circle}). A setup must keep the three
crossings within the detector's span, or bring the outliers in through its top or bottom face.
Figure~\ref{fig:curve} shows the reference's crossing height over a wider range of fields than the
level's three; where it lies above the detector's span, as without field, the flight may still
come down into the top face, so the figure alone does not decide which fields the reference
tolerates. A map of the fields or placements that solve would take full flights or the game's
verdicts, which this document does not reproduce.

\section{What the game computes}\label{sec:engine}

The game adds the disturbance's uniform $B_z$ to the charges' fields (\physref{2.3}) and integrates
\eqref{eq:eom-b} with its adaptive eighth-order Runge--Kutta method (\physref{5}, \physref{6},
\physref{7}), each disturbance's flight verified on its own. The problem is the same as here, and
the differences in the comparison are the two integrations' errors.

\section{Further reading}\label{sec:reading}

Motion in a uniform magnetic field: \Jackson{§12.2}, \levelref{first_coil}. Flights under
disturbances: \levelref{stray_field}, \levelref{mains_hum}.

\levelfileappendix

\end{document}
